% JCIM-style review manuscript: single column with wide, readable floats.
\documentclass[manuscript,nonacm]{acmart}

\AtBeginDocument{\providecommand\BibTeX{{Bib\TeX}}}
\settopmatter{printacmref=false}
\renewcommand{\footnotetextcopyrightpermission}[1]{}
\pagestyle{plain}
\raggedbottom
\setlength{\emergencystretch}{4em}

\usepackage{booktabs}
\usepackage{amsmath}
\usepackage{graphicx}
\usepackage{multirow}
\usepackage{array}
\usepackage{longtable}
\usepackage{xspace}
\usepackage{microtype}
\usepackage{placeins}
\usepackage{float}
\usepackage{pdflscape}

\newcommand{\method}{TAPER\xspace}
\newcommand{\mace}{MoleculeACE\xspace}

\begin{document}

\title{Target-Excluded Response Memory for Generalizable Activity-Cliff Prediction}

\author{Anonymous Authors}
\affiliation{\institution{Anonymous Institution}\country{}}
\email{anonymous@example.com}
\renewcommand{\shortauthors}{Anonymous et al.}

\begin{abstract}
Activity cliffs violate the local smoothness assumption underlying many
structure-based activity predictors: nearly identical molecules can have
substantially different potency against the same target. We hypothesize that a
molecule's measurements in other biological contexts provide functional
information complementary to its chemical structure. We introduce \method as a
target-excluded response-memory framework: each molecule is represented by
collision-free chemical environments and by its ChEMBL response profile after
removing the prediction target and all biologically equivalent targets. Sparse
coverage, standardized potency, assay-resolved context, and fit-selected dense
memory coordinates are combined with chemistry in a single
positive-semidefinite kernel and fitted by target-specific support vector
regression. On 30 MoleculeACE targets, three repeated training-only five-fold
evaluations reduce macro RMSE by $8.01\%$ overall and $9.32\%$ on cliff
molecules relative to the identical chemistry-only predictor. The frozen model
confirms on the official test split. Component restrictions, sensitivity
analyses, and molecule-level cases on the same benchmark show that the gain
depends on molecule-specific, target-excluded response geometry rather than
nearest-label transfer or indiscriminate feature addition. External regression
and matched-pair classification collections show that the same response-memory
principle transfers beyond the primary MoleculeACE regression setting.
\end{abstract}

\keywords{activity cliffs, response memory, bioactivity profiles, target exclusion, ChEMBL, kernel regression, quantitative structure--activity relationships}

\maketitle

\section{Introduction}
Activity cliffs (ACs) are pairs of structurally similar compounds with large
potency differences. They occupy the least smooth regions of a
structure--activity landscape: a model may achieve a low average error while
systematically smoothing away the local changes that medicinal chemists care
about. \mace formalizes this problem with 30 target-specific regression
datasets and separate errors for cliff and non-cliff molecules~\cite{moleculeace}.
Graph encoders, molecular pretraining, and cliff-aware objectives improve the
structural learner, but they continue to infer target-specific response mainly
from ligand structure~\cite{scage,graphcliff,acanets}.

Our starting point is that chemical identity and biological response answer
different questions. Structure describes what a molecule is. Measurements
against other targets and in other assays describe how that molecule has
behaved across biological contexts. Such measurements cannot determine the
current endpoint directly, but they can distinguish two chemically similar
molecules whose functional profiles are different. This distinction is most
valuable near a cliff, where structural proximity is least informative.

Using external profiles creates an equally important statistical problem. A
record for the benchmark target, or for the same protein under another name,
would leak the answer. Sparse assay coverage can also turn chance correlations
into apparently predictive features. We therefore formulate the method around
target exclusion rather than unrestricted database reuse. For every benchmark
target, the current protein and all biologically equivalent ChEMBL entities are
removed before any profile is constructed. Profile coordinates,
normalization, relevance selection, and kernel scales are then estimated only
from the fitting fold.

The resulting method, \method (Target-Excluded Activity Profile--Enhanced Regression),
combines a collision-free multi-view chemical kernel with target- and
assay-resolved activity-profile kernels. It is a single supervised regressor,
not a cascade of a base model and a post-hoc correction. It neither averages
neighbor labels nor stores model residuals. Our contributions are: (i) a
target-excluded activity-profile formulation with explicit biological-identity
control; (ii) a positive-semidefinite chemistry--profile kernel that represents
both observation coverage and relative potency; (iii) repeated train-only
evaluation, component restrictions, sensitivity analyses, and held-out case
analyses on \mace that directly test the proposed mechanism. Low-coverage
GraphCliff assays and ACNet matched-pair tasks are transfer checks that test
the same response-memory principle under additional endpoints.

\section{Related Work}
The benchmark foundation is \mace, which shows that extended-connectivity
fingerprints can outperform many graph models in low-data and cliff-heavy
regimes~\cite{moleculeace}. ACNet instead frames cliff recognition as
matched-molecular-pair classification~\cite{acnet}. SCAGE combines
self-conformation-aware pretraining, functional-group tasks, and attention-based
substructure interpretation~\cite{scage}. GraphCliff uses short--long-range
gating to reduce graph oversmoothing~\cite{graphcliff}, whereas ACANet couples
regression to a cliff-aware soft-margin triplet objective~\cite{acanets}. These
methods improve how structure or structural differences are encoded.

Bioactivity profiling provides a complementary representation: compounds are
compared by patterns of measured response rather than by structure alone. Its
use for held-out prediction, however, depends on assay overlap and explicit
leakage control. Recent work shows that profile-based cliff prediction is
sensitive to assay-neighborhood construction and that apparently transferable
chemical transformations weaken after leakage is removed~\cite{bioactivityprofiles,kinasecliffs}.
PrismNet provides a broader structural comparison by decomposing scaffolds,
functional groups, and pharmacophores across 64 datasets~\cite{prismnet}.
\method differs from both lines: it keeps the strong local inductive bias of
chemical fingerprints, but conditions chemical similarity on a profile from
which the evaluated biological endpoint has been explicitly removed.

\section{Method}
\subsection{Problem formulation and design principle}
For each benchmark target $t$, the training set is
$\mathcal{D}_t=\{(s_i,y_i)\}_{i=1}^{N_t}$, where $s_i$ is a canonical SMILES
and $y_i$ is the measured $pK_i$ or $pEC_{50}$. Cliff annotations are used only
to stratify the reported error; the proposed model never consumes them as
features, weights, or auxiliary labels. We learn one target-specific function
$f_t$ from a precomputed kernel. The design follows one hypothesis: two
molecules are stronger analogues when they are similar in both chemical
structure and target-excluded biological response. Every chemical vocabulary,
profile statistic, relevance score, and kernel scale is estimated from the
current fitting fold. Validation and test molecules are projected only through
those frozen transformations.

\subsection{Target-excluded activity profiles}
Let the external activity archive be
$\mathcal{A}=\{(m_k,e_k,a_k,v_k)\}_{k=1}^{N_{\mathcal A}}$, where $m_k$
identifies a standardized molecule, $e_k$ a biological target entity, $a_k$ an
assay, and $v_k$ a standardized potency measurement. For a benchmark target
$t$, let $\mathcal{B}_t$ contain $t$ and every ChEMBL entity judged
biologically equivalent to it. The admissible activity profile of a query
molecule $q$ is
\begin{equation}
\mathcal{P}_t(q)=
\{(e_k,a_k,v_k):(m_k=q)\land(e_k\notin\mathcal{B}_t)\}.
\end{equation}
The equality $m_k=q$ denotes exact standardized molecular identity; the method
does not fetch labels from chemically similar compounds. We exclude entities
sharing a UniProt accession, gene-symbol synonym, or normalized preferred name
with the benchmark target. Repeated molecule--entity and molecule--assay
measurements are aggregated by their median. If $\mathcal{P}_t(q)$ is empty,
all profile coordinates are zero and prediction reduces to the chemical branch.

\begin{figure*}[t]
\centering
\includegraphics[width=0.98\textwidth]{figures/framework-cliff.png}
\Description{Workflow of target-excluded response memory. A molecule's ChEMBL bioactivity profile is filtered to remove the prediction target, then combined with multi-view molecular representations and unified kernel regression to improve cliff-aware prediction.}
\caption{Method overview. \method builds a target-excluded response memory by
removing the prediction target and biologically equivalent targets from each
molecule's external bioactivity profile. The resulting response vector is fused
with multi-view chemical representations through a unified kernel and fitted by
target-specific SVR. The final prediction uses response memory to distinguish
structurally similar molecules that can have different activity near cliffs.}
\label{fig:method}
\end{figure*}

\subsection{Collision-free multi-view chemistry}
The chemistry branch deliberately retains the strong local-inductive bias of
fingerprints. Let $\mathcal{G}$ denote a fixed collection of complementary
environment generators spanning circular, feature-aware, and atom-pair
descriptions at multiple resolutions. For each generator
$g\in\mathcal{G}$, sparse integer environment identifiers are collected on
the fitting molecules, avoiding collisions caused by folding environments into
a fixed bit vector.  We construct both a binary vector $b_i^g$ and a count
vector $c_i^g$.  With generalized Tanimoto similarity
\begin{equation}
T(x,x')=\frac{x^\top x'}{\lVert x\rVert_2^2+\lVert x'\rVert_2^2-x^\top x'},
\end{equation}
the multi-view chemical kernel is
\begin{equation}
K_{\mathrm{chem}}(i,i')=\frac{1}{2|\mathcal{G}|}\sum_{g\in\mathcal{G}}
\left[T(b_i^g,b_{i'}^g)+T(c_i^g,c_{i'}^g)\right].
\end{equation}
An unseen environment in a held-out molecule has no fitting-vocabulary column,
which is equivalent to an out-of-vocabulary zero and cannot alter training
features.

\subsection{Activity-profile representation}
\paragraph{Coverage and potency geometry.}
For profile feature $j$, $v_{ij}$ denotes molecule $i$'s external potency,
and fitting-fold observations define its robust location $\mu_j$ and scale
$\sigma_j$. We store an observation indicator
$o_{ij}=\mathbb{1}[v_{ij}\ \text{observed}]$ and a centered response
$z_{ij}=o_{ij}(v_{ij}-\mu_j)/\max(\sigma_j,\varepsilon_{\rm scale})$, where
$\varepsilon_{\rm scale}>0$ prevents division by a vanishing scale. After row-wise
$\ell_2$ normalization, a basic response kernel is
\begin{equation}
U(i,i')=\frac{\langle \bar{o}_i,\bar{o}_{i'}\rangle+
\langle \bar{z}_i,\bar{z}_{i'}\rangle}{2}.
\end{equation}
The first term distinguishes shared experimental coverage from missingness;
the second measures agreement in relative potency. We build $U_{\rm ent}$ at
the target-entity level and $U_{\rm assay}$ at the assay level. Entity features
pool measurements that address the same molecular target; assay features keep
experimental contexts separate. Their complementarity is tested directly by
removing either family from the final kernel.

\paragraph{Fit-only relevance and dense views.}
Many off-target columns are irrelevant. For every feature whose fitting support
$n_j$ exceeds a prespecified minimum $n_{\min}$, we compute
\begin{equation}
w_j=\max\left(\operatorname{corr}(v_{\cdot j},y)^2-
\frac{1}{n_j-1},0\right).
\end{equation}
Here $w_j$ is a finite-sample-corrected relevance score, and the correlation is
computed only where feature $j$ is observed. The subtraction is the analytic
null expectation of squared correlation and suppresses small-support
coincidences. The $q_{\rm resp}$ highest positive $w_j$ values select
target-level and assay-level dense views. Each selected view concatenates
standardized potency and its observation indicator, then applies an RBF kernel.
Its feature location, scale, and RBF bandwidth---the median positive fitting-pair
squared distance---are all fitting-fold quantities.  Denote the resulting
kernels by $D_{\rm ent}$ and $D_{\rm assay}$.

Let $\mathcal{V}_{\rm prof}$ collect the sparse entity, sparse assay, dense
entity, dense assay, and two chemistry--response interaction views. The
complete activity-profile block is
\begin{align}
K_{\mathrm{prof}}=\frac{1}{|\mathcal{V}_{\rm prof}|}\big(&U_{\rm ent}+U_{\rm assay}
+D_{\rm ent}+D_{\rm assay}\nonumber\\
&+K_{\rm chem}\odot D_{\rm ent}
+K_{\rm chem}\odot D_{\rm assay}\big),
\end{align}
where $\odot$ is the Hadamard product.  The two interaction terms require
molecules to agree simultaneously in local chemistry and response behavior.

\subsection{Chemistry--profile fusion and model fitting}
Let $\mathcal{R}=\{\mathrm{chem},\mathrm{prof}\}$ index the two retained
positive-semidefinite evidence blocks. They are fused as
\begin{equation}
K_{\boldsymbol{\lambda}}(i,i')=
\sum_{r\in\mathcal{R}}\lambda_r K_r(i,i'),
\qquad \lambda_r\geq0,\quad \sum_{r\in\mathcal{R}}\lambda_r=1.
\end{equation}
Here $K_r$ is the kernel contributed by representation block $r$, and $\lambda_r$ is
its global mixture coefficient. The coefficients are fixed once using only the
training-side component audit and are then shared by every target, split, and query;
they are not target-adaptive quantities. Non-negative sums preserve positive
semidefiniteness, and the Schur product theorem covers the Hadamard
interactions. On $K_{\boldsymbol{\lambda}}$, a single
$\epsilon$-insensitive support vector regressor solves
\begin{equation}
\min_{f\in\mathcal{H}_{K_{\boldsymbol{\lambda}}}}\frac{1}{2}
\lVert f\rVert_{\mathcal{H}_{K_{\boldsymbol{\lambda}}}}^{2}
+C\sum_i\max\big(|y_i-f(s_i)|-\epsilon,0\big),
\end{equation}
where $\mathcal{H}_{K_{\boldsymbol{\lambda}}}$ is the reproducing-kernel Hilbert
space, $C$ controls regularization, and $\epsilon$ defines the insensitive
regression tube. Their fixed values and all preprocessing settings are reported
in the experimental specification rather than encoded in the model equation.
Training proceeds independently for each target and fitting fold. We first
construct the collision-free chemical vocabulary, then compute the admissible
profiles $\mathcal{P}_t$ and all fit-only profile statistics. The resulting
training--training kernel is passed to one $\epsilon$-insensitive support vector
regressor; there is no neural pretraining, epoch selection, or post-hoc
correction model. At evaluation, we construct only the query--training
cross-kernel with the frozen transformations and apply the learned support-vector
coefficients. There is no per-target tuning, model ensemble, cliff reweighting,
or validation/test-label feedback in the frozen model. The prediction is
\begin{equation}
\widehat y_q=b+\sum_{i\in\mathcal S}\alpha_i
K_{\boldsymbol{\lambda}}\!\left(
(q,\mathcal P_t(q)),(s_i,\mathcal P_t(s_i))\right),
\end{equation}
where $\mathcal S$ is the fitted support-vector set, $\alpha_i$ its learned
coefficient for molecule $i$, and $b$ the intercept. Chemical similarity is
contained in $K_{\boldsymbol{\lambda}}$ even though the compact notation shows
the activity profile explicitly. The external archive supplies covariates;
supervised endpoint activities enter only through the fitted regressor. Thus
\method is one profile-conditioned predictor rather than a chemistry model
followed by a residual or nearest-neighbor correction.

For a molecule-level audit, we rank fitting molecules by the chemistry block
$K_{\rm chem}(q,i)$. This ranking visualizes the structural neighborhood seen
by the fitted kernel; no ranked label is copied or averaged during prediction.
It allows us to ask whether activity profiles correct cases in which chemical
neighbors alone imply the wrong local trend.

\paragraph{Falsifiable mechanism.}
The activity-profile hypothesis makes four testable predictions. First,
adding a target-excluded profile should improve the identical chemistry
predictor under repeated held-out splits. Second, entity-level response,
assay-level context, sparse coverage, dense potency, and chemistry--profile
agreement should carry complementary information. Third, the gain should
disappear when the same profile rows are randomly reassigned to other
molecules, because that preserves dimensionality and coverage while removing
molecule-specific biological information. Fourth, improvement should include
cases in which the closest structural labels are misleading; otherwise the
method would be only a more elaborate expression of chemical similarity. The
Results section tests these predictions primarily within \mace, using repeated
splits, restricted representations, profile controls, sensitivity analyses, and
molecule-level cases.

\begin{table*}[t]
\centering
\caption{Dataset statistics. A supervised unit is one activity record for the regression collections and one matched molecular pair for ACNet. The event column denotes cliff molecules for regression and positive pairs for ACNet; percentages are relative to supervised units.}
\label{tab:datasets}
\small
\begin{tabular}{lrrrrrrr}
\toprule
Collection & Targets & Unique mol. & Units & Events (\%) & Train & Validation & Test\\
\midrule
MoleculeACE & 30 & 35,633 & 48,714 & 18,831 (38.7) & 38,912 & -- & 9,802\\
GraphCliff external & 5 & 296 & 296 & 122 (41.2) & 234 & -- & 62\\
ACNet & 190 & 53,085 & 402,079 & 20,861 (5.2) & 321,590 & 40,119 & 40,370\\
\bottomrule
\end{tabular}
\end{table*}

\begin{table*}[t]
\centering
\caption{Task-separated target-macro comparisons. Panel A contains only locally evaluated MoleculeACE methods. Panel B re-runs every Panel-A architecture on the five GraphCliff external assays under the same train-only five-fold protocol, and additionally reports the vendored MoleculeACE AFP, CNN, GAT, GCN, and MPNN architectures. Panel C evaluates a matched-pair \method{}-edit instantiation on ACNet. Mean $\mathbin{\pm}$ target-level standard deviation is reported when per-target values are available. Lower RMSE is better; higher ROC--AUC and AP are better.}
\label{tab:main}
\scriptsize
\renewcommand{\arraystretch}{1.02}
\begin{tabular*}{0.86\textwidth}{@{\extracolsep{\fill}}lrr}
\toprule
\multicolumn{3}{l}{\textit{A. MoleculeACE regression (30 targets)}}\\
Method & RMSE & Cliff RMSE\\
\midrule
Linear ECFP & 0.835 $\mathbin{\pm}$ 0.109 & 0.867 $\mathbin{\pm}$ 0.080\\
ExtraTrees ECFP & 0.703 $\mathbin{\pm}$ 0.075 & 0.798 $\mathbin{\pm}$ 0.062\\
Random forest ECFP & 0.731 $\mathbin{\pm}$ 0.080 & 0.816 $\mathbin{\pm}$ 0.066\\
ECFP-SVM & 0.673 $\mathbin{\pm}$ 0.070 & 0.765 $\mathbin{\pm}$ 0.057\\
HistGBM ECFP & 0.721 $\mathbin{\pm}$ 0.077 & 0.808 $\mathbin{\pm}$ 0.058\\
MLP ECFP & 0.855 $\mathbin{\pm}$ 0.102 & 0.880 $\mathbin{\pm}$ 0.075\\
Chemprop & 0.882 $\mathbin{\pm}$ 0.096 & 0.941 $\mathbin{\pm}$ 0.088\\
GINE + NodeNorm & 0.901 $\mathbin{\pm}$ 0.108 & 0.958 $\mathbin{\pm}$ 0.098\\
GraphCliff & 0.858 $\mathbin{\pm}$ 0.102 & 0.928 $\mathbin{\pm}$ 0.092\\
\textbf{\method} & \textbf{0.569 $\mathbin{\pm}$ 0.078} & \textbf{0.638 $\mathbin{\pm}$ 0.092}\\
\midrule
\multicolumn{3}{l}{\textit{B. GraphCliff external regression (5 assays)}}\\
Method & RMSE & Cliff RMSE\\
\midrule
Linear ECFP & 0.588 $\mathbin{\pm}$ 0.117 & 0.685 $\mathbin{\pm}$ 0.258\\
RBF-SVM ECFP & 0.618 $\mathbin{\pm}$ 0.162 & 0.718 $\mathbin{\pm}$ 0.310\\
HistGBM ECFP & 1.046 $\mathbin{\pm}$ 0.401 & 1.091 $\mathbin{\pm}$ 0.662\\
MLP ECFP & 0.946 $\mathbin{\pm}$ 0.247 & 0.911 $\mathbin{\pm}$ 0.351\\
ExtraTrees ECFP & 0.713 $\mathbin{\pm}$ 0.146 & 0.817 $\mathbin{\pm}$ 0.273\\
Random forest ECFP & 0.761 $\mathbin{\pm}$ 0.172 & 0.837 $\mathbin{\pm}$ 0.295\\
Chemprop & 0.782 $\mathbin{\pm}$ 0.141 & 0.801 $\mathbin{\pm}$ 0.259\\
GINE + NodeNorm & 0.795 $\mathbin{\pm}$ 0.161 & 0.790 $\mathbin{\pm}$ 0.240\\
GraphCliff encoder & 0.974 $\mathbin{\pm}$ 0.339 & 1.056 $\mathbin{\pm}$ 0.460\\
AFP & 1.394 $\mathbin{\pm}$ 0.182 & 1.375 $\mathbin{\pm}$ 0.195\\
CNN & 0.970 $\mathbin{\pm}$ 0.213 & 1.055 $\mathbin{\pm}$ 0.490\\
GAT & 1.419 $\mathbin{\pm}$ 0.341 & 1.573 $\mathbin{\pm}$ 0.595\\
GCN & 1.438 $\mathbin{\pm}$ 0.286 & 1.542 $\mathbin{\pm}$ 0.568\\
MPNN & 1.393 $\mathbin{\pm}$ 0.333 & 1.445 $\mathbin{\pm}$ 0.716\\
\textbf{\method} & \textbf{0.530 $\mathbin{\pm}$ 0.156} & \textbf{0.596 $\mathbin{\pm}$ 0.300}\\
\midrule
\multicolumn{3}{l}{\textit{C. ACNet pair classification (190 targets)}}\\
Method & ROC--AUC & AP\\
\midrule
Linear ECFP & 0.864 $\mathbin{\pm}$ 0.144 & 0.407 $\mathbin{\pm}$ 0.274\\
RBF-SVM ECFP & 0.927 $\mathbin{\pm}$ 0.109 & 0.695 $\mathbin{\pm}$ 0.275\\
HistGBM ECFP & 0.910 $\mathbin{\pm}$ 0.148 & 0.678 $\mathbin{\pm}$ 0.278\\
MLP ECFP & 0.785 $\mathbin{\pm}$ 0.238 & 0.470 $\mathbin{\pm}$ 0.316\\
ExtraTrees ECFP & 0.933 $\mathbin{\pm}$ 0.112 & 0.696 $\mathbin{\pm}$ 0.270\\
Random forest ECFP & 0.929 $\mathbin{\pm}$ 0.116 & 0.678 $\mathbin{\pm}$ 0.270\\
Chemprop + pair head & 0.782 $\mathbin{\pm}$ 0.166 & 0.397 $\mathbin{\pm}$ 0.312\\
GINE + PairNorm + pair head & 0.769 $\mathbin{\pm}$ 0.171 & 0.351 $\mathbin{\pm}$ 0.282\\
GraphCliff encoder + pair head & 0.710 $\mathbin{\pm}$ 0.223 & 0.329 $\mathbin{\pm}$ 0.301\\
\textbf{\method{}-edit} & \textbf{0.949 $\mathbin{\pm}$ 0.091} & \textbf{0.753 $\mathbin{\pm}$ 0.261}\\
\bottomrule
\end{tabular*}
\end{table*}


\section{Experimental Protocol}
\paragraph{Primary regression benchmark.}
We use all 30 official \mace targets and preserve the byte-level train/test
split. Model specification is performed only inside the official training pool.
We first compare restricted activity-profile representations and directly remove the
chemistry--response interaction under a stratified five-fold split. Candidate
branches without paired evidence are discarded, after which we repeat the
frozen minimal model under two additional split seeds. Every profile statistic,
relevance score, chemical vocabulary,
and kernel scale is re-estimated inside the fitting fold. Uncertainty is
computed by resampling targets rather than treating folds as independent.

\paragraph{Profile identity and leakage controls.}
The response-profile shuffle control keeps the same molecule set, selected
profile coordinates, feature dimensionality, coverage distribution, profile
value distribution, and kernel construction. Within each fitting fold and its
validation partition, complete profile rows are randomly reassigned across
molecules before the profile kernels are recomputed and fused with the
unchanged chemical kernel. The only destroyed information is therefore the
identity link between a molecule and its own external biological-response
profile. For the leakage stress test, we deliberately re-inject the current
benchmark endpoint into the external archive for the official-training
molecules, then compare strict biological exclusion, normalized target-name
exclusion, and no exclusion under the same folds. The no-exclusion condition is
intentionally contaminated and is used only to quantify how much apparent
performance can be created by endpoint leakage.

\paragraph{Supplementary transfer checks.}
The paper is centered on \mace. We additionally report five low-coverage
GraphCliff assays as an external stress test for profile transfer and an ACNet
matched-pair classification appendix as an orthogonal edit-information check.
These secondary collections do not define the method, do not share the primary
metric, and are not pooled with the \mace regression evidence.

\paragraph{Confirmation and audit status.}
After the train-only specification is frozen, the final artifact is refit on
all official training rows and evaluated on the official test partition. Since
earlier component audits had already inspected this test resource, we do not call it
a pristine lockbox; the repeated train-only CV remains the primary
generalization evidence. Crucially, official-test results are not used to add,
remove, or reweight a component.

\paragraph{Post-fit case selection.}
Case studies are illustrative and selected only after the model is frozen;
they are not estimates of a representative effect. We define sparse, moderate,
and dense response coverage as one to two, three to seven, and at least eight
observed response features. Eligible official-test queries are cliff molecules
with at most 30 heavy atoms, positive error reduction, and nearest-chemistry
similarity of at least $0.65$. Within each stratum we select the largest error
reduction and use dataset identifier and SMILES as deterministic tie breaks.
The two fitting molecules with highest $K_{\rm chem}$ are then displayed.
Test labels are used only for this post-fit selection and visualization.

The reproduced official ECFP-SVM is the primary historical baseline, and the
collision-free binary+count kernel is the strongest same-regressor, same-split
chemistry control. Table~\ref{tab:main} is restricted to the main \mace
regression benchmark. It includes locally refitted ECFP baselines, Chemprop,
the strongest local GINE variant, and GraphCliff. Every architecture retained
in this primary panel is rerun under the same external GraphCliff protocol;
literature-only ImageMol, GEM, SCAGE, and MTPNet values are not mixed into the
local comparison. The vendored MoleculeACE AFP, CNN, GAT, GCN, and MPNN
architectures are additionally run on the external assays. The
per-target Table~\ref{tab:moleculeace_all} provides the detailed local
five-fold comparison for every locally run MoleculeACE method. The like-for-like
local baseline audit is summarized separately in
Table~\ref{tab:local-cv-significance} and
Figure~\ref{fig:local-baseline-significance}.

\begin{table*}[t]
\centering
\caption{Local five-fold RMSE on all 30 MoleculeACE targets. Every method in this table is evaluated locally; fold results are averaged within each target before ranking. Lower is better; bold and underline denote the best and second-best method in each row. Reduction is relative to the strongest local baseline excluding \method.}
\label{tab:moleculeace_all}
\tiny
\setlength{\tabcolsep}{3.0pt}
\renewcommand{\arraystretch}{1.02}
\resizebox{\textwidth}{!}{%
\begin{tabular}{lrrrrrrrrrrrrr}
\toprule
Target & CV $n$ & Linear & RF & ExtraTrees & SVM & HGB & MLP & Chemprop & GINE-Node & GraphCliff & MTPNet & \method & Reduction (\%) $\uparrow$\\
\midrule
CHEMBL1862\_Ki & 1,899 & 0.880 & 0.769 & 0.741 & \underline{0.710} & 0.762 & 0.888 & 0.929 & 0.951 & 0.942 & 0.764 & \textbf{0.660} & +7.1\\
CHEMBL1871\_Ki & 1,575 & 0.725 & 0.644 & 0.630 & \underline{0.623} & 0.659 & 0.907 & 0.752 & 0.743 & 0.762 & 0.660 & \textbf{0.590} & +5.3\\
CHEMBL2034\_Ki & 1,794 & 0.722 & 0.702 & 0.669 & \underline{0.645} & 0.685 & 0.821 & 0.909 & 0.852 & 0.795 & 0.738 & \textbf{0.610} & +5.3\\
CHEMBL2047\_EC50 & 1,509 & 0.706 & 0.660 & 0.634 & \textbf{0.616} & 0.650 & 0.784 & 0.871 & 0.864 & 0.803 & 0.704 & \underline{0.618} & -0.2\\
CHEMBL204\_Ki & 6,603 & 0.965 & 0.838 & 0.798 & \underline{0.723} & 0.807 & 0.813 & 1.013 & 1.052 & 0.984 & 0.816 & \textbf{0.648} & +10.4\\
CHEMBL2147\_Ki & 3,486 & 0.854 & 0.762 & 0.720 & \underline{0.669} & 0.758 & 0.857 & 1.023 & 1.070 & 0.905 & 0.686 & \textbf{0.538} & +19.5\\
CHEMBL214\_Ki & 7,953 & 0.821 & 0.702 & 0.672 & \underline{0.649} & 0.689 & 0.771 & 0.803 & 0.873 & 0.780 & 0.756 & \textbf{0.583} & +10.2\\
CHEMBL218\_EC50 & 2,469 & 0.819 & 0.735 & 0.706 & \underline{0.655} & 0.698 & 0.842 & 0.898 & 0.911 & 0.844 & 0.741 & \textbf{0.617} & +5.8\\
CHEMBL219\_Ki & 4,473 & 0.879 & 0.765 & 0.741 & \underline{0.717} & 0.757 & 0.906 & 0.907 & 0.927 & 0.872 & 0.747 & \textbf{0.628} & +12.3\\
CHEMBL228\_Ki & 4,086 & 0.918 & 0.808 & 0.774 & 0.732 & 0.792 & 0.929 & 0.954 & 0.961 & 0.930 & \textbf{0.690} & \underline{0.699} & -1.3\\
CHEMBL231\_Ki & 2,328 & 0.835 & 0.731 & 0.702 & \underline{0.673} & 0.714 & 0.937 & 0.834 & 0.847 & 0.817 & 0.835 & \textbf{0.577} & +14.3\\
CHEMBL233\_Ki & 7,536 & 0.924 & 0.772 & 0.744 & \underline{0.721} & 0.757 & 0.878 & 0.851 & 0.908 & 0.830 & 0.802 & \textbf{0.588} & +18.4\\
CHEMBL234\_Ki & 8,772 & 0.832 & 0.710 & 0.680 & \underline{0.648} & 0.695 & 0.799 & 0.820 & 0.877 & 0.804 & 0.682 & \textbf{0.554} & +14.4\\
CHEMBL235\_EC50 & 5,631 & 0.863 & 0.711 & 0.683 & \underline{0.668} & 0.708 & 0.780 & 0.840 & 0.886 & 0.845 & 0.679 & \textbf{0.618} & +7.5\\
CHEMBL236\_Ki & 6,231 & 0.855 & 0.726 & 0.695 & \underline{0.672} & 0.720 & 0.795 & 0.863 & 0.910 & 0.840 & 0.719 & \textbf{0.578} & +14.0\\
CHEMBL237\_EC50 & 2,286 & 1.067 & 0.916 & 0.879 & \underline{0.845} & 0.923 & 1.023 & 1.119 & 1.151 & 1.156 & 0.879 & \textbf{0.724} & +14.3\\
CHEMBL237\_Ki & 6,243 & 0.965 & 0.759 & 0.729 & \underline{0.723} & 0.747 & 0.906 & 0.889 & 0.925 & 0.849 & 0.745 & \textbf{0.596} & +17.5\\
CHEMBL238\_Ki & 2,517 & 0.816 & 0.713 & 0.687 & \underline{0.661} & 0.699 & 1.009 & 0.832 & 0.834 & 0.788 & 0.748 & \textbf{0.582} & +12.0\\
CHEMBL239\_EC50 & 4,122 & 0.874 & 0.745 & 0.722 & \underline{0.701} & 0.732 & 0.801 & 0.900 & 0.920 & 0.871 & 0.734 & \textbf{0.645} & +7.9\\
CHEMBL244\_Ki & 7,428 & 0.948 & 0.830 & 0.790 & \underline{0.730} & 0.813 & 0.870 & 1.010 & 1.117 & 1.024 & 0.816 & \textbf{0.660} & +9.6\\
CHEMBL262\_Ki & 2,049 & 0.947 & 0.764 & 0.738 & 0.716 & 0.766 & 0.970 & 0.931 & 0.928 & 1.003 & \underline{0.707} & \textbf{0.607} & +14.1\\
CHEMBL264\_Ki & 6,864 & 0.800 & 0.705 & 0.678 & \underline{0.636} & 0.691 & 0.799 & 0.820 & 0.875 & 0.776 & 0.643 & \textbf{0.588} & +7.5\\
CHEMBL2835\_Ki & 1,467 & 0.473 & 0.436 & 0.425 & \underline{0.411} & 0.441 & 0.584 & 0.569 & 0.556 & 0.561 & 0.419 & \textbf{0.311} & +24.2\\
CHEMBL287\_Ki & 3,183 & 0.881 & 0.799 & 0.770 & \underline{0.716} & 0.772 & 1.062 & 0.886 & 0.879 & 0.875 & 0.815 & \textbf{0.626} & +12.6\\
CHEMBL2971\_Ki & 2,337 & 0.792 & 0.684 & 0.665 & 0.641 & 0.687 & 0.864 & 0.919 & 0.890 & 0.855 & \underline{0.585} & \textbf{0.503} & +14.1\\
CHEMBL3979\_EC50 & 2,697 & 0.747 & 0.689 & 0.667 & \underline{0.644} & 0.695 & 0.756 & 0.818 & 0.830 & 0.807 & 0.660 & \textbf{0.589} & +8.5\\
CHEMBL4005\_Ki & 2,301 & 0.700 & 0.656 & 0.628 & \underline{0.598} & 0.649 & 0.703 & 0.873 & 0.874 & 0.836 & 0.674 & \textbf{0.573} & +4.2\\
CHEMBL4203\_Ki & 1,746 & 0.866 & 0.752 & 0.727 & \underline{0.707} & 0.745 & 1.014 & 0.939 & 0.925 & 0.909 & 0.857 & \textbf{0.514} & +27.3\\
CHEMBL4616\_EC50 & 1,629 & 0.823 & 0.710 & \underline{0.699} & 0.706 & 0.727 & 0.794 & 0.841 & 0.820 & 0.813 & 0.759 & \textbf{0.675} & +3.5\\
CHEMBL4792\_Ki & 3,522 & 0.757 & 0.732 & 0.705 & \underline{0.638} & 0.706 & 0.802 & 0.862 & 0.880 & 0.860 & 0.785 & \textbf{0.481} & +24.6\\
\midrule
Macro mean & 116,736 & 0.835 & 0.731 & 0.703 & \underline{0.673} & 0.721 & 0.855 & 0.882 & 0.901 & 0.858 & 0.728 & \textbf{0.593} & +11.9\\
Best count & -- & 0 & 0 & 0 & 1 & 0 & 0 & 0 & 0 & 0 & 1 & \textbf{28} & --\\
Mean rank & -- & 8.17 & 5.37 & 3.27 & 2.10 & 4.63 & 8.47 & 9.53 & 10.17 & 8.47 & 4.77 & \textbf{1.07} & --\\
\bottomrule
\end{tabular}
}
\end{table*}


\section{Results}
\subsection{Target-excluded profiles improve the same chemical predictor}
Across three independently shuffled five-fold splits, the chemistry baseline
has macro RMSE $0.6428$, cliff RMSE $0.7380$, and non-cliff RMSE $0.5854$.
\method reaches $0.5913$, $0.6693$, and $0.5437$, reductions of $8.01\%$,
$9.32\%$, and $7.13\%$. The per-seed overall reductions are $7.97\%$,
$7.99\%$, and $8.08\%$; the conclusion is therefore insensitive to the split
seed. Across all 450 paired folds, target-clustered
bootstrap intervals for the absolute reduction are $[0.0372,0.0678]$ overall
and $[0.0465,0.0953]$ on cliffs. Target-averaged paired tests give
$p<2\times10^{-8}$ for both endpoints. The improvement is positive for 29/30
targets overall and 27/30 targets on cliffs after averaging seeds.

The official-test confirmation gives RMSE $0.5687$ overall, $0.6377$ on cliffs,
and $0.5266$ on non-cliffs. Against the reproduced official ECFP-SVM, these are
reductions of $15.3\%$, $15.1\%$, and $15.6\%$. Because the resource had been
inspected in earlier component audits, we treat this as confirmation and base the
central claim on repeated training-only CV.

Under the local five-fold audit in Table~\ref{tab:moleculeace_all}, \method has
lower overall RMSE than the strongest locally evaluated baseline on 28 of 30
targets after averaging folds within each target. The mean paired gain against
this per-target best local baseline is $0.0766$ RMSE. Relative to the strongest
baseline by macro average, ECFP-SVM ($0.6733$), \method reaches $0.5927$, an
$11.9\%$ reduction. The two local exceptions are CHEMBL2047\_EC50 and
CHEMBL228\_Ki.

\begin{table}[t]
\centering
\caption{Local five-fold MoleculeACE comparison. Values are mean $\mathbin{\pm}$ standard deviation across 30 targets after averaging folds within each target. $P$ values are Holm-adjusted one-sided paired Wilcoxon tests of whether each baseline has larger target-level error than \method.}
\label{tab:local-cv-significance}
\scriptsize
\begin{tabular}{lrrr}
\toprule
Method & RMSE & Cliff RMSE & $P$ (cliff)\\
\midrule
Linear ECFP & 0.834 $\mathbin{\pm}$ 0.110 & 0.871 $\mathbin{\pm}$ 0.082 & $<0.001$\\
RF-ECFP & 0.731 $\mathbin{\pm}$ 0.082 & 0.817 $\mathbin{\pm}$ 0.064 & $<0.001$\\
ExtraTrees-ECFP & 0.704 $\mathbin{\pm}$ 0.076 & 0.801 $\mathbin{\pm}$ 0.059 & $<0.001$\\
ECFP-SVM & 0.673 $\mathbin{\pm}$ 0.070 & 0.768 $\mathbin{\pm}$ 0.056 & $<0.001$\\
HGB-ECFP & 0.723 $\mathbin{\pm}$ 0.079 & 0.812 $\mathbin{\pm}$ 0.058 & $<0.001$\\
MLP-ECFP & 0.858 $\mathbin{\pm}$ 0.100 & 0.882 $\mathbin{\pm}$ 0.076 & $<0.001$\\
Chemprop & 0.882 $\mathbin{\pm}$ 0.096 & 0.941 $\mathbin{\pm}$ 0.088 & $<0.001$\\
GINE + NodeNorm & 0.901 $\mathbin{\pm}$ 0.108 & 0.958 $\mathbin{\pm}$ 0.098 & $<0.001$\\
GraphCliff & 0.858 $\mathbin{\pm}$ 0.102 & 0.928 $\mathbin{\pm}$ 0.092 & $<0.001$\\
MTPNet & 0.728 $\mathbin{\pm}$ 0.089 & 0.821 $\mathbin{\pm}$ 0.133 & $<0.001$\\
\textbf{\method} & \textbf{0.593 $\mathbin{\pm}$ 0.076} & \textbf{0.673 $\mathbin{\pm}$ 0.069} & --\\
\bottomrule
\end{tabular}
\end{table}


\begin{figure*}[t]
\centering
\includegraphics[width=0.96\textwidth]{figures/fig_local_baseline_significance.pdf}
\Description{Horizontal bars compare locally fitted baselines and TAPER on overall and cliff RMSE. Error bars show target-level standard deviations and annotations report paired target-level significance against TAPER.}
\caption{Like-for-like local five-fold comparison on all 30 MoleculeACE
targets. Bars and error bars denote the target mean and standard deviation
after folds are averaged within each target. Annotations give one-sided paired
Wilcoxon $P$ values for the hypothesis that the corresponding baseline has
higher target-level error than \method. Values are Holm-adjusted across the
displayed baselines, and folds are not treated as independent replicates.}
\label{fig:local-baseline-significance}
\end{figure*}

The fully local audit asks a stricter question: after all baselines are run
under the same five-fold protocol, does \method still beat the best baseline
available for each target rather than only a fixed reference method? It does.
Relative to the per-target best local baseline, \method reduces target-mean
RMSE by $0.0766$ overall, $0.0699$ on cliff molecules, and $0.0701$ on
non-cliff molecules. The corresponding win counts are 28/30, 27/30, and 27/30.
The profile-evidence analysis is intentionally conservative: the number of
positive response-profile coordinates has a positive association with
cliff-RMSE gain ($\rho=0.34$), but this remains a trend rather than a
standalone claim ($p=0.066$). The stronger conclusion comes from the paired
gain distribution and the component-level restrictions below.

\begin{figure*}[t]
\centering
\includegraphics[width=0.96\textwidth]{figures/fig_baseline_advantage_analysis.pdf}
\Description{Panel a shows target-level RMSE gains of TAPER over the strongest local baseline for overall, cliff, and non-cliff molecules. Panel b summarizes mean paired gains with target-bootstrap intervals and target win counts. Panel c relates cliff gain to the number of positive response-profile coordinates.}
\caption{Where the local advantage over baselines comes from. For every target
and endpoint, the comparison baseline is the strongest locally refit baseline
on that target after folds are averaged. Panel (a) shows paired target-level
gain distributions. Panel (b) reports the mean paired gain with 95\%
target-bootstrap intervals and win counts. Panel (c) tests whether targets with
more positive response-profile coordinates tend to obtain larger cliff gains;
the positive association is suggestive but not significant at the conventional
5\% level.}
\label{fig:baseline-advantage}
\end{figure*}

\begin{figure}[t]
\centering
\includegraphics[width=0.92\linewidth]{figures/fig_inference_efficiency.pdf}
\Description{Scatter plot of locally evaluated MoleculeACE methods with overall RMSE on the x-axis and cliff RMSE on the y-axis. Each method has a distinct marker shape and is directly labeled. TAPER lies in the lower-left region.}
\caption{Local baseline error trade-off on MoleculeACE. Each point is one
locally evaluated method after folds are averaged within targets and then
macro-averaged over the 30 targets. The horizontal axis shows overall RMSE and
the vertical axis shows cliff RMSE; the lower-left region is better. Marker
shape distinguishes methods, and direct labels replace method-family legends.
Only locally run methods are included.}
\label{fig:efficiency}
\end{figure}

The local scatter makes the baseline relationship explicit. The strongest
classical baseline is ECFP-SVM, whereas the locally refitted graph baselines are
substantially less competitive under this low-data cliff benchmark. \method is
the only method in the lower-left region, showing simultaneous improvement on
overall and cliff RMSE rather than a trade-off between the two endpoints.

\subsection{Entity, assay, coverage, and potency views are complementary}
Table~\ref{tab:ablation} and Figure~\ref{fig:ablation} restrict the
activity-profile representation under one five-fold split while keeping the
regressor and chemistry branch unchanged. Entity aggregation pools measurements
that address the same biological target; assay aggregation preserves
experimental context. Sparse views compare observed coverage and relative
response directly, whereas dense views first select fitting-fold coordinates
that carry reproducible endpoint information. Chemistry combined with either
entity or assay profiles is weaker than their joint representation; likewise,
dense or sparse response geometry alone is weaker than their combination.
Restricting the external profile to protein-sequence neighbors changes overall
RMSE from $0.5928$ to $0.5927$ and therefore earns no place in the final model.
The literal leave-one-block-out comparison removes chemistry--profile
interactions; doing so
increases overall RMSE by $0.0032$ (25/30 target-level increases; one-sided
Wilcoxon $p=2.5\times10^{-5}$) and cliff RMSE by $0.0021$ (18/30 increases;
$p=0.060$). Thus the interaction has reproducible support on the overall
endpoint, while its cliff-specific effect remains uncertain. The larger
restricted-view gaps show that no single profile geometry substitutes for
the complete representation.

\begin{table}[t]
\centering
\caption{Activity-profile representations on 30 targets (five-fold CV). Lower
is better. Every row uses the same regressor and chemistry branch; the first
row omits external activity profiles.}
\label{tab:ablation}
\footnotesize
\begin{tabular}{lccc}
\toprule
Retrieval strategy & Overall & Cliff & Non-cliff\\
\midrule
Chemistry only & 0.6440 & 0.7417 & 0.5855\\
Dense profiles only & 0.6127 & 0.6959 & 0.5615\\
Assay-level profiles only & 0.6052 & 0.6884 & 0.5552\\
Entity-level profiles only & 0.6032 & 0.6853 & 0.5529\\
Sparse profiles only & 0.6072 & 0.6904 & 0.5583\\
No chem.--profile interaction & 0.5959 & 0.6749 & 0.5482\\
\textbf{Hybrid sparse+dense, entity+assay} & \textbf{0.5927} & \textbf{0.6727} & \textbf{0.5443}\\
\bottomrule
\end{tabular}
\end{table}

\begin{figure*}[t]
\centering
\includegraphics[width=0.82\textwidth]{figures/fig_component_ablation.pdf}
\Description{Horizontal box plots show target-level overall and cliff RMSE penalties for chemistry alone, four restricted chemistry-plus-profile representations, and one interaction removal, each compared with complete TAPER.}
\caption{Penalty of restricted representations relative to the complete model
(30 target-level paired differences). Boxes show medians and interquartile
ranges, whiskers extend to $1.5$ times the interquartile range, and translucent
points show individual targets. The combined response representation
outperforms each restricted view. Removing the chemistry--response interaction
gives a small overall penalty; its cliff-specific target-level test remains
inconclusive.}
\label{fig:ablation}
\end{figure*}

The same audit removes attractive but unsupported complexity. Protein-sequence
nearest-neighbor views do not improve any endpoint; a local physical kernel has
a sub-percent, non-significant overall effect ($p=0.13$); and the susceptibility
extension does not significantly improve cliff error. These branches are not
part of \method and remain documented only as negative controls.

\subsection{Supplementary transfer check: low-coverage external assays}
On five newly standardized, low-coverage GraphCliff assays, the unchanged model
reduces macro RMSE from $0.5558$ to $0.5299$ ($4.66\%$) and cliff RMSE from
$0.6445$ to $0.5957$ ($7.58\%$). Non-cliff RMSE worsens by $3.31\%$, identifying
sparse response coverage as a real boundary rather than hiding it in the
average. Among the locally refit graph encoders, NodeNorm GINE is strongest
overall on these assays (RMSE $0.798$), but remains behind both the linear ECFP
control and \method; this is consistent with the very small per-assay fitting
pools rather than evidence that a larger graph encoder is automatically useful.

As a second cliff-focused stress test, we locally ran two-fold evaluation on
the dopamine D2, factor Xa, and SARS-CoV-2 Mpro collections released by
Dablander et al. A calibrated matched-molecular-pair correction improves the
ECFP-SVM on every dataset: the cross-dataset overall RMSE changes from
$0.6982$ to $0.6896$, while cliff RMSE changes from $0.9787$ to $0.9584$.
This experiment supports transferable local edit information, but is kept
separate from the activity-profile claim because it neither uses target-excluded
profiles nor evaluates the complete \method model.

\begin{table}[t]
\centering
\caption{Local two-fold evaluation on three external DABlander activity-cliff datasets. Values are mean $\mathbin{\pm}$ standard deviation across datasets after averaging folds within each dataset. Lower is better.}
\label{tab:dablander-external}
\small
\begin{tabular}{lrr}
\toprule
Method & RMSE & Cliff RMSE\\
\midrule
ECFP-SVM & 0.698 $\mathbin{\pm}$ 0.092 & 0.979 $\mathbin{\pm}$ 0.045\\
Calibrated MMP & 0.690 $\mathbin{\pm}$ 0.092 & 0.958 $\mathbin{\pm}$ 0.049\\
\bottomrule
\end{tabular}
\end{table}


\begin{table}[t]
\centering
\caption{Five-fold RMSE on every external GraphCliff assay. The baseline is the collision-free chemistry kernel using the same regressor and folds.}
\label{tab:external_all}
\scriptsize
\begin{tabular}{lrrrr}
\toprule
Assay & $n$ & Chemistry & Ours & $\Delta$\%\\
\midrule
CHEMBL2311243\_IC50 & 51 & 0.491 & 0.495 & -0.8\\
CHEMBL2598\_IC50 & 48 & 0.771 & 0.726 & +5.8\\
CHEMBL3024\_IC50 & 73 & 0.348 & 0.297 & +14.9\\
CHEMBL4685\_IC50 & 78 & 0.653 & 0.581 & +11.0\\
CHEMBL5014\_IC50 & 46 & 0.516 & 0.551 & -6.8\\
\bottomrule
\end{tabular}
\end{table}


\subsection{Response-memory transfer to matched-pair classification}
ACNet tests the same central claim in a pair-classification form: activity
cliffs are better modeled when the predictor represents not only two molecular
endpoints, but also the local response induced by the edit between them. We
therefore instantiate a task-specific \method{}-edit model by combining endpoint
fingerprints with an explicit chemical-edit memory. On all 190 ACNet target
tasks, this \method{}-edit variant raises target-macro
ROC--AUC from $0.9334$ to $0.9490$, average precision from $0.6961$ to $0.7529$,
and accuracy from $0.9345$ to $0.9483$. Target-bootstrap 95\% intervals for the
paired gains exclude zero for ROC--AUC ($[0.0094,0.0225]$), average precision
($[0.0406,0.0736]$), and accuracy ($[0.0089,0.0187]$); the corresponding
win/tie/loss counts are $103/64/23$, $104/66/20$, and $112/64/14$. This result
supports the broader response-memory hypothesis under a different endpoint:
local biological or edit responses provide information that endpoint structure
alone misses. The strongest refit graph-pair baseline, PairNorm GINE, reaches
ROC--AUC $0.769$ and average precision $0.351$, showing that the ACNet gain is
not explained by replacing endpoint fingerprints with a graph encoder.

\begin{table}[t]
\centering
\caption{Locally executed ACNet comparison over the same 190 targets. Values are target-macro mean $\mathbin{\pm}$ standard deviation. Higher is better.}
\label{tab:acnet-local-macro}
\scriptsize
\begin{tabular}{lrr}
\toprule
Method & ROC--AUC & AP\\
\midrule
Linear ECFP & 0.864 $\mathbin{\pm}$ 0.144 & 0.407 $\mathbin{\pm}$ 0.274\\
RF-ECFP & 0.929 $\mathbin{\pm}$ 0.116 & 0.678 $\mathbin{\pm}$ 0.270\\
ExtraTrees-ECFP & 0.933 $\mathbin{\pm}$ 0.112 & 0.696 $\mathbin{\pm}$ 0.270\\
ECFP-SVM & 0.927 $\mathbin{\pm}$ 0.109 & 0.695 $\mathbin{\pm}$ 0.275\\
HGB-ECFP & 0.910 $\mathbin{\pm}$ 0.148 & 0.678 $\mathbin{\pm}$ 0.278\\
MLP-ECFP & 0.785 $\mathbin{\pm}$ 0.238 & 0.470 $\mathbin{\pm}$ 0.316\\
Chemprop & 0.782 $\mathbin{\pm}$ 0.166 & 0.397 $\mathbin{\pm}$ 0.312\\
GINE + PairNorm & 0.769 $\mathbin{\pm}$ 0.171 & 0.351 $\mathbin{\pm}$ 0.282\\
GraphCliff & 0.710 $\mathbin{\pm}$ 0.223 & 0.329 $\mathbin{\pm}$ 0.301\\
\bottomrule
\end{tabular}
\end{table}


\subsection{Sensitivity identifies the useful profile regime}
Figure~\ref{fig:sensitivity} varies one internal parameter at a time while
retaining the complete model. Expanding the fit-selected dense-response budget
from one to ten coordinates reduces cliff RMSE from $0.6825$ to
$0.6727$ and overall RMSE from $0.5965$ to $0.5927$, with the latter saturating
after eight coordinates. Thus several complementary response coordinates are
more useful than one selected assay. Performance is nearly invariant to a
four-fold contraction or expansion of the median-derived dense-kernel
bandwidth. The profile mixture has a broad useful region: intermediate weights
remain near the minimum, whereas overweighting the external profile degrades
both endpoints. Finally, a small SVR penalty underfits sharply; performance
improves up to the fixed setting and then plateaus. Thus profile dimensionality,
evidence balance, and predictor regularization materially affect the result,
while dense-view bandwidth is comparatively insensitive.

\begin{figure*}[t]
\centering
\includegraphics[width=0.82\textwidth]{figures/fig_sensitivity.pdf}
\Description{Four line plots show overall and cliff RMSE as the response-coordinate budget, dense-kernel bandwidth, profile mixture weight, and SVR penalty are varied one at a time.}
\caption{One-parameter-at-a-time sensitivity of the complete \method model.
Panels vary (a) the number of fit-selected response coordinates, (b) the
dense-response RBF bandwidth relative to its fitting-fold median, (c) the
profile coefficient in chemistry--profile fusion, and (d) the SVR
penalty. Points are means over 30 target-level five-fold averages; shaded bands
are 95\% target-bootstrap intervals. Orange rings and vertical dashed lines mark
the fixed manuscript setting. Every model component remains active throughout;
within each panel, only the named parameter changes.}
\label{fig:sensitivity}
\end{figure*}

\FloatBarrier

\subsection{Cold-start coverage, retrieval controls, and discovery ranking}
We first replaced the released random-like partition with one fixed 80/20
Bemis--Murcko scaffold split for each of all 30 targets. No scaffold is shared
between fitting and held-out molecules. The held-out scaffold uses only the
archived, target-excluded profile records available under this protocol; thus a
molecule without such a record receives zero response memory rather than a
profile inferred from its structural neighbors. Under this stricter split,
the locally refit GINE+NodeNorm graph baseline has macro overall/cliff RMSE
$0.876/0.968$ and the chemistry-only kernel has $0.732/0.819$, while \method
reaches $0.701/0.769$ (a $4.3\%/6.1\%$ reduction from the chemistry kernel).
The absolute performance drop relative to the released split is expected, but
the retained paired improvement shows that the response memory is not
restricted to the original molecular split.

\begin{figure*}[t]
\centering
\includegraphics[width=0.80\textwidth]{figures/fig_scaffold_coldstart.pdf}
\Description{Paired target-level overall and cliff RMSE values for ECFP-SVM, locally refit GINE+NodeNorm, and TAPER under 30 independently scaffold-disjoint splits.}
\caption{Scaffold-disjoint cold-start evaluation. Each line is one target under
a fixed 80/20 Bemis--Murcko split; points are target-level errors and black
segments join macro means. The test scaffold cannot obtain a profile by
structural-neighbor imputation.}
\label{fig:scaffold-coldstart}
\end{figure*}

We next stress-tested the response-memory claim on the frozen official-test
predictions rather than changing the model after seeing the labels.  Splitting
test molecules by the number of external response records gives a clear
cold-start boundary: with no record, overall RMSE is $0.600$ versus $0.650$
for ECFP-SVM, whereas the cliff subset is $0.710$ versus $0.789$.  The gain
persists with one record ($0.061$ mean absolute-error reduction) and two to
three records ($0.069$), and reaches $0.167$ above ten records (target-cluster
bootstrap intervals are stored with the analysis).  These are post-hoc
coverage strata, not a new split or a claim that coverage is randomly
assigned; they quantify the expected cold-start cost of exact-identity memory.

The result is not explained by a single nearest-response lookup.  Among the
7,164 test molecules with at least one response record, the mean absolute-error
gain remains positive in every quartile of nearest-response similarity
($0.063$--$0.110$).  The profile-only quantities are therefore useful
diagnostics of evidence availability, while the final prediction still uses
the full chemistry--profile kernel and all training molecules.

As a formal control, we also ran the profile-only and late-fusion kernels on
all 30 targets with the same fit-only vocabulary and biological exclusion.
The response-only kernel is weak (macro RMSE $0.969$, cliff RMSE $1.038$),
whereas adding response coverage to chemistry gives $0.628/0.715$ and adding
the potency-residual profile gives $0.620/0.702$ (overall/cliff RMSE).  The
combined off-target profile fusion reaches $0.609/0.692$, demonstrating that
the gain comes from a complementary profile--chemistry representation rather
than from retrieval alone.

We made the retrieval control explicit on the same 30 target-held-out splits.
Profile-$k$NN retrieves the five closest response profiles and averages their
training activities; profile kernel ridge regression (KRR) uses only the
target-excluded response-profile kernel; and a fixed equal late fusion averages
the ECFP-SVM and profile-KRR predictions.  These profile-only controls are
substantially weaker (macro RMSE $1.049$ for $k$NN and $1.134$ for KRR) than
ECFP-SVM ($0.640$); fixed late fusion is also worse ($0.776$).  This controls
the alternative explanation that performance is caused by a nearest-profile
lookup or by appending an otherwise generic profile predictor.  In contrast,
the jointly constructed response kernel improves the chemistry kernel
($0.609$ versus $0.640$).

\begin{figure*}[t]
\centering
\includegraphics[width=0.78\textwidth]{figures/fig_profile_retrieval_controls.pdf}
\Description{Target-level overall and cliff RMSE distributions for ECFP-SVM, profile k-nearest-neighbor retrieval, profile kernel ridge regression, and their fixed prediction-level late fusion.}
\caption{Explicit retrieval controls. Each point is one target under the same
official split and target-exclusion protocol. Profile-only $k$NN and KRR, and
a fixed late fusion of ECFP-SVM with KRR, do not reproduce the performance of
the integrated response kernel. Horizontal bars are target-macro means.}
\label{fig:profile-retrieval-controls}
\end{figure*}

\subsection{Causal dose, temporal, and biological-radius controls}
Natural response coverage is confounded by compound history: molecules with
many recorded assays need not be equally difficult.  We therefore performed a
within-molecule dose experiment.  For each target, we retained only the same
official-test molecules with at least 16 target-excluded assay records, froze
the training-side kernel and SVR, and randomly kept $0,1,2,4,8$, or 16 records
per molecule for 20 independent masks.  Structures, prediction targets, and
test molecules do not change.  Across the 15 targets with sufficient rich
cliff/non-cliff subsets, both overall and cliff performance improve
monotonically as records are restored: the target-macro overall/cliff RMSE
reduction relative to chemistry changes from $-0.199/-0.246$ with no record
to $0.015/0.005$ with 16 records.  This is a causal manipulation of memory
availability, rather than an observational coverage stratification.  The
restricted rich cohort does not show a positive plateau by 16 records; we
therefore interpret it as evidence for a dose-dependent recovery, not a claim
that memory is already saturated.

The previous manipulation freezes the fitted model and changes only query-side
memory.  We therefore performed a stricter matched train/test experiment in
which both partitions were restricted to the same high-coverage molecules and
were independently masked at every dose before rebuilding the kernel and
refitting the SVR.  Labels, structures, molecule identities, and the official
split were fixed.  Across 13 targets that retained both cliff classes at all
doses, the reduction is numerically zero at dose zero ($-0.000/-0.001$ for
overall/cliff RMSE) and increases to $0.066/0.068$ at 16 records per molecule
(20 masks per dose).  Thus the dose response is not an artifact of keeping a
fully observed training representation while only manipulating test queries.

We also separated biological potency from the history of which assays were
observed.  Under strict exclusion on all 30 targets, we held every
molecule--response-coordinate observation mask fixed and deterministically
permuted only observed potency values within each coordinate.  This preserves
coverage, dimensionality, coordinate-wise value distributions, and the full
kernel construction, while destroying molecule-specific biological responses.
The true profiles achieve $0.568/0.638$ overall/cliff RMSE, whereas the
fixed-mask potency shuffle gives $0.620/0.707$.  The loss after shuffling
therefore rules out an explanation based solely on assay-history width or an
additional feature channel.

We next asked whether the profile merely contains a near-duplicate target.
On all 30 released test targets, we successively removed target-equivalent
records, sequences with at least 80\% identity, sequences with at least 50\%
identity, and all targets sharing an InterPro/Pfam/PANTHER protein-family
identifier.  Overall/cliff RMSE remains $0.568/0.638$ after the 80\% radius,
changes to $0.589/0.671$ after the 50\% radius, and is $0.623/0.719$ after
the family-level exclusion.  Thus the strongest signal is not explained by an
80\%-identical proxy, although the family-level loss shows that related
biology is genuinely useful information rather than irrelevant extra width.

That family-level degradation could still arise merely because family exclusion
removes many more records.  For each target we consequently created ten
seeded random target-exclusion sets on top of strict exclusion, selecting
whole response coordinates to match the family-only deletion count.  The
entity-plus-assay record-count mismatch is at most one record (mean 0.003).
Count-matched random exclusion yields $0.603/0.690$ overall/cliff RMSE,
whereas the actual family exclusion yields $0.623/0.719$.  Related biological
responses are therefore specifically informative: deleting the same quantity
of unrelated response history is materially less harmful.

The same-external-data-budget controls make this distinction explicit.  Every
baseline receives the same ECFP, response-observation mask, and standardized
response vector; the full vector is compressed by a label-free fit-only random
projection for the RBF-SVM.  Concatenated RBF-SVM reaches $0.625/0.673$ and
train-only tuned late fusion reaches $0.634/0.707$, whereas the integrated
response kernel reaches $0.568/0.638$ (overall/cliff RMSE).  Hence response
memory itself contains information, but joint chemistry--response geometry
extracts more of it than concatenation or a tuned prediction-level blend.

Finally, we constructed a dated ChEMBL audit on the predeclared CHEMBL244
target.  For 3,097 molecules with a benchmark document year, each response
record was retained only if its document year was no later than that molecule's
benchmark measurement.  On the 621 held-out dated molecules, chemistry-only
RMSE is $0.648$; the complete archive and pre-measurement profiles give
$0.624$ and $0.623$, respectively (cliff RMSE $0.709$ and $0.711$ versus
$0.746$).  This single-target date-verified subset cannot establish a general
prospective claim, but it directly rules out a purely future-document
explanation for this observed transfer gain.

\begin{figure*}[t]
\centering
\includegraphics[width=0.94\textwidth]{figures/fig_memory_robustness.pdf}
\Description{A six-panel figure shows matched train-test response-record dose response, a fixed observation-mask potency shuffle, family versus count-matched random biological exclusion, equal-data-budget baselines, widening biological exclusion radius, and a document-year temporal profile cutoff.}
\caption{Robustness controls for the response-memory claim. (a) The same rich
train and test molecules are independently masked at every dose before the
kernel is rebuilt and the SVR is refit; intervals are target-bootstrap 95\%
intervals over 13 matched targets. (b) Potency values are shuffled within a
response coordinate while every molecule-coordinate observation is retained.
(c) Family exclusion is compared with whole-target random exclusions that
match the deleted response-record count. (d) Each competing baseline receives
the same response records, masks, and standardized response values; points are
30 target-level cliff errors. (e) Wider sequence and family exclusion removes
possible target proxies. (f) A dated CHEMBL244 subset uses only response
documents available no later than each molecule's benchmark measurement.}
\label{fig:memory-robustness}
\end{figure*}

Finally, we evaluated a fixed early-discovery ranking protocol on every target:
the top $10\%$ of observed test potency defines hits, and scores are ranked
without using those labels.  The locally trained Chemprop comparator has macro
EF1\%/ROC--AUC/BEDROC-like values of $7.33/0.888/0.259$; ECFP-SVM gives
$8.12/0.910/0.283$, and \method gives $8.75/0.935/0.313$.  Improvements are
not universal at the individual-target level, so these values are reported as
a discovery-oriented retrospective ranking check rather than as a claim of a
prospective screen.

We also performed a strict ChEMBL-to-BindingDB transfer test.  For each
MoleculeACE target, the regressor was fit only on its released ChEMBL training
partition and evaluated on BindingDB molecules absent from that partition.
All BindingDB systems biologically equivalent to the prediction protein,
including the systems supplying the external labels, were removed before the
response profile was constructed.  Across 26 targets with sufficient external
data (7,941 unseen molecules), the frozen $10{:}8$ chemistry--response kernel
reduces macro RMSE from $0.928$ to $0.916$ and cliff RMSE from $1.248$ to
$1.221$.  Only $17.8\%$ of external molecules have an observed retained
response coordinate, so this is a deliberately sparse, no-recalibration
transfer test rather than a database-specific re-tuning exercise.

\begin{figure*}[t]
\centering
\includegraphics[width=0.94\textwidth]{figures/fig_extended_evidence.pdf}
\Description{Four panels show response-memory coverage scaling, gain by nearest-response similarity quartile, virtual-screening enrichment, and the relation between response-coordinate availability and external cross-validation error.}
\caption{Extended evidence from frozen local runs. (a) Overall and cliff RMSE
by exact-identity response coverage; (b) positive absolute-error gain across
nearest-response similarity quartiles, showing that the method is not a single
nearest-neighbor label transfer; (c) target-macro early-discovery ranking
metrics; and (d) external cross-validation RMSE as a function of the number of
fit-derived response coordinates. All statistics are computed from stored
predictions or training-fold summaries; no test label is used for model
selection.}
\label{fig:extended-evidence}
\end{figure*}

\begin{table*}[t]
\centering
\small
\caption{Deterministically selected held-out case-study extremes. Positive gain favours ResponseKernel.}
\begin{tabular}{llrrr}
\toprule
Target & Case & $y$ & $|e|$ gain & Response records\\
\midrule
CHEMBL1862\_Ki & largest penalty & 6.70 & -0.80 & 2\\
CHEMBL1862\_Ki & largest rescue & 4.10 & +1.27 & 43\\
CHEMBL1871\_Ki & largest penalty & 5.13 & -0.49 & 3\\
CHEMBL1871\_Ki & largest rescue & 7.36 & +0.57 & 0\\
CHEMBL2034\_Ki & largest penalty & 5.90 & -0.53 & 1\\
CHEMBL2034\_Ki & largest rescue & 8.89 & +0.93 & 1\\
CHEMBL2047\_EC50 & largest penalty & 8.11 & -0.61 & 1\\
CHEMBL2047\_EC50 & largest rescue & 4.98 & +0.74 & 0\\
CHEMBL204\_Ki & largest penalty & 5.77 & -0.95 & 0\\
CHEMBL204\_Ki & largest rescue & 10.84 & +1.70 & 4\\
CHEMBL2147\_Ki & largest penalty & 8.40 & -1.64 & 0\\
CHEMBL2147\_Ki & largest rescue & 7.30 & +1.31 & 38\\
\bottomrule
\end{tabular}
\label{tab:extended-cases}
\end{table*}


\FloatBarrier

\subsection{What biological response dimensions are selected?}
To make the response memory biologically auditable, we annotated the
fit-selected assay coordinates using the cached ChEMBL preferred names across
the two repeated-CV runs. Receptor-associated coordinates account for 1,092
selection events, kinases for 428, enzymes for 112, ion channels for 75, and
transporters for 51; 891 events remain in a deliberately broad ``other''
category that includes non-protein and incompletely annotated ChEMBL targets.
The most recurrent individual coordinates include the mu opioid receptor,
5-hydroxytryptamine receptors 1E and 1D, PPAR$\gamma$, CYP2C9/CYP2D6, KCNH2,
and CAMK1D. Thus the memory is pharmacological rather than a single assay
modality. This frequency analysis describes what the fit-only relevance
procedure selects; it is not a causal statement that any one off-target
explains a potency measurement.

\begin{figure*}[t]
\centering
\includegraphics[width=0.80\textwidth]{figures/fig_response_biology.pdf}
\Description{Bar plots show pharmacological-family counts and the most frequently selected ChEMBL response coordinates across repeated target-fold runs.}
\caption{Biological annotation of fit-selected response coordinates. (a) Counts
by conservative preferred-name category; (b) the most frequently selected
individual ChEMBL target coordinates. Selection is performed within the fitting
fold after prediction-target exclusion.}
\label{fig:response-biology}
\end{figure*}

\FloatBarrier

\subsection{Activity profiles resolve misleading structural neighborhoods}
Response coverage is a candidate explanation, not a universal guarantee. Across
the 30 targets, the association between the number of informative response
features and cliff-RMSE reduction is positive but inconclusive
($\rho=0.25$, bootstrap 95\% CI $[-0.18,0.63]$, $p=0.18$). Across a complete
leave-one-target-out sensitivity analysis, $\rho$ ranges from $0.17$ to $0.36$.
The largest-gain target is CHEMBL4203, but the observations do not justify a
fitted coverage--gain trend.

Across all 9,802 official-test molecules, the median absolute-error reduction
is $0.0506$; among 3,790 cliff molecules it is $0.0601$. The displayed partner
examples are selected after model freezing and are not representative effect
estimates. We select the largest positive cliff rescue with at most 30 heavy
atoms in each of three predeclared response-coverage strata, then display the
two most similar official-training molecules under $K_{\rm chem}$.
Separately, the machine-readable extended-case corpus contains the largest
rescue and largest penalty for every one of the 30 targets (60 held-out
molecules in total), along with structure, response-record count, structural
and response similarity, and all three local predictions. Table~\ref{tab:extended-cases}
shows a compact deterministic subset; the full corpus is retained without
post-hoc filtering.

\begin{figure*}[t]
\centering
\includegraphics[width=0.98\textwidth]{figures/fig_retrieval_cases.pdf}
\Description{Three held-out activity-cliff queries, each beside its two closest official-training molecules under the frozen chemistry kernel. A monochrome dot-line comparison gives absolute error for ECFP-SVM and TAPER.}
\caption{Post-fit structural-neighborhood cases. Each row shows one held-out
query, its two closest official-training molecules under frozen
$K_{\rm chem}$, and the change in absolute error from ECFP-SVM to \method.
Molecules are aligned to the query and rendered in monochrome. The displayed
neighbors explain the chemistry-only context; their labels are not transferred
by \method.}
\label{fig:retrieval-cases}
\end{figure*}

The cases expose why response context helps. For CHEMBL4792, the two closest
training analogues have activities $7.28$ and $6.30$, while the difluoro query is
$5.30$; the smooth ECFP-SVM stays high at $6.91$, whereas \method predicts
$5.41$. For CHEMBL219, neighbor activities $5.82$ and $6.02$ bracket the query
at $5.96$, but ECFP-SVM predicts $7.38$ and \method $5.84$. The sharpest
counterexample to label transfer is CHEMBL262: both structural analogues are
potent ($8.30$ and $7.80$), yet the query is $6.00$. Chemistry alone predicts
$7.28$; the dense activity-profile view moves the estimate to $6.03$. Thus the
gain is not nearest-label copying: chemistry locates the local neighborhood,
while cross-context response geometry determines whether its apparent
smoothness should be trusted.

\FloatBarrier

\section{Discussion and Limitations}
The central result supports a specific representation principle: chemical
structure and target-excluded activity profiles contain complementary
information near activity cliffs. Physical descriptors, sequence-selected
profile neighborhoods, and an explicit susceptibility channel do not earn
their complexity and are excluded from the final model. This negative evidence
matters because the gain cannot be attributed to indiscriminate feature
addition. Equation~(6) therefore uses only a global chemistry--profile mixture,
shared across targets, splits, and queries, with strict biological exclusion
applied before every profile is constructed.

Five limitations matter. First, external-response availability is not random;
well-studied compounds and targets have denser profiles. Second, the BindingDB
transfer covers 26 exact-protein target sets but is still retrospective and
does not establish prospective assay performance. Third, SCAGE, GraphCliff, ACANet, and PrismNet do not all
publish identical splits or aggregations, so their table entries are contextual
rather than the basis of our central claim.
Fourth, our method uses additional ChEMBL measurements; a fully controlled
leaderboard should equalize external-data budgets. The 190 ACNet tasks increase
breadth only for the edit-representation hypothesis and must not be pooled with
regression RMSE. Fifth, exact-identity profile assignment means that an unseen
compound with no ChEMBL record receives no external response information; the
present model does not infer a profile from similar molecules. The
scaffold-disjoint evaluation shows a retained advantage under this conservative
condition, but does not remove the coverage limitation or replace a prospective
new-compound study.

\section{Conclusion}
\method tests whether activity measured outside the prediction target can make
chemical similarity more informative without leaking the endpoint. Across
repeated MoleculeACE splits, target-excluded entity and assay profiles improve
the same collision-free chemical regressor, with the largest absolute reduction
on cliff molecules. Shuffle controls show that the gain requires
molecule-specific response identity, and deliberate leakage controls show why
strict biological exclusion is required. The MoleculeACE
restricted-representation analyses and held-out cases support one explanation:
structure locates a chemical neighborhood, whereas cross-context response
distinguishes molecules that the neighborhood would otherwise smooth together.
The method remains dependent on external profile availability and does not solve
cliffs for compounds absent from the archive. These boundaries make
target-excluded activity profiling a testable MoleculeACE-grounded
representation principle rather than a universal substitute for structural or
interaction-based modeling.

\bibliographystyle{ACM-Reference-Format}
\bibliography{references}

\clearpage
\onecolumn
\appendix
\begin{landscape}
\section{All ACNet target results}
Table~\ref{tab:acnet_all} reports the held-out ROC--AUC for every matched-pair
task. It is included to expose the full positive, tied, and negative target
distribution rather than only the macro average.
\fontsize{6.2}{6.35}\selectfont
\setlength{\tabcolsep}{3.0pt}
\begin{longtable}{@{}rrrrr@{\hspace{1.5em}}rrrrr@{\hspace{1.5em}}rrrrr@{}}
\caption{Test ROC--AUC for all 190 ACNet target tasks using the ExtraTrees head selected globally by validation macro average precision. Endpoint fingerprints form the baseline; Ours adds the explicit edit representation.}\label{tab:acnet_all}\\
\toprule
Target & Pairs & Endpoint & Ours & $\Delta$ & Target & Pairs & Endpoint & Ours & $\Delta$ & Target & Pairs & Endpoint & Ours & $\Delta$\\
\midrule
\endfirsthead
\multicolumn{15}{c}{\tablename\ \thetable{} -- continued}\\
\toprule
Target & Pairs & Endpoint & Ours & $\Delta$ & Target & Pairs & Endpoint & Ours & $\Delta$ & Target & Pairs & Endpoint & Ours & $\Delta$\\
\midrule
\endhead
\midrule\multicolumn{15}{r}{Continued on next page}\\\endfoot
\bottomrule\endlastfoot
8 & 4,861 & 0.989 & 0.994 & +0.005 & 252 & 14,073 & 0.971 & 0.975 & +0.004 & 11451 & 670 & 0.955 & 0.970 & +0.015\\
11 & 2,952 & 0.977 & 0.969 & -0.008 & 259 & 9,041 & 0.948 & 0.962 & +0.013 & 11473 & 200 & 1.000 & 1.000 & +0.000\\
15 & 5,393 & 0.909 & 0.904 & -0.004 & 262 & 439 & 0.883 & 0.911 & +0.029 & 11522 & 582 & 0.974 & 0.991 & +0.018\\
28 & 149 & 1.000 & 1.000 & +0.000 & 278 & 5,491 & 0.979 & 0.986 & +0.007 & 11534 & 1,474 & 0.871 & 0.941 & +0.070\\
36 & 432 & 1.000 & 1.000 & +0.000 & 280 & 10,253 & 0.979 & 0.987 & +0.008 & 11536 & 768 & 0.990 & 1.000 & +0.010\\
43 & 677 & 0.701 & 0.806 & +0.104 & 10009 & 5,371 & 0.993 & 0.996 & +0.003 & 11635 & 515 & 0.912 & 0.837 & -0.075\\
47 & 528 & 1.000 & 1.000 & +0.000 & 10034 & 540 & 0.943 & 0.943 & +0.000 & 11682 & 109 & 1.000 & 1.000 & +0.000\\
48 & 55 & 0.800 & 1.000 & +0.200 & 10044 & 124 & 1.000 & 1.000 & +0.000 & 11848 & 136 & 0.926 & 0.981 & +0.056\\
50 & 576 & 0.956 & 0.991 & +0.035 & 10102 & 23,704 & 0.986 & 0.993 & +0.006 & 11869 & 303 & 0.988 & 1.000 & +0.012\\
51 & 5,747 & 0.954 & 0.977 & +0.023 & 10131 & 403 & 0.975 & 0.975 & +0.000 & 11871 & 337 & 0.939 & 0.909 & -0.030\\
52 & 692 & 1.000 & 1.000 & +0.000 & 10139 & 50 & 1.000 & 1.000 & +0.000 & 11942 & 658 & 0.939 & 0.970 & +0.030\\
56 & 943 & 0.394 & 0.649 & +0.255 & 10142 & 8,020 & 0.983 & 0.979 & -0.004 & 12000 & 1,049 & 0.982 & 0.984 & +0.002\\
59 & 153 & 1.000 & 1.000 & +0.000 & 10167 & 114 & 1.000 & 1.000 & +0.000 & 12054 & 191 & 1.000 & 1.000 & +0.000\\
65 & 633 & 0.976 & 0.976 & +0.000 & 10184 & 1,268 & 0.992 & 0.992 & +0.000 & 12186 & 635 & 0.968 & 0.968 & +0.000\\
69 & 127 & 0.950 & 0.875 & -0.075 & 10188 & 481 & 0.906 & 0.922 & +0.017 & 12209 & 4,331 & 0.946 & 0.953 & +0.007\\
72 & 26,376 & 0.980 & 0.967 & -0.014 & 10193 & 5,229 & 0.932 & 0.943 & +0.011 & 12252 & 1,064 & 0.934 & 0.915 & -0.019\\
86 & 670 & 0.911 & 0.958 & +0.047 & 10197 & 193 & 0.980 & 1.000 & +0.020 & 12592 & 421 & 0.869 & 0.960 & +0.091\\
87 & 7,043 & 0.931 & 0.934 & +0.003 & 10209 & 2,361 & 0.951 & 0.962 & +0.011 & 12659 & 619 & 0.290 & 0.339 & +0.048\\
88 & 1,973 & 1.000 & 1.000 & +0.000 & 10244 & 347 & 0.971 & 0.971 & +0.000 & 12666 & 59 & 1.000 & 1.000 & +0.000\\
90 & 1,923 & 0.969 & 0.942 & -0.026 & 10260 & 10,782 & 0.992 & 0.994 & +0.003 & 12679 & 79 & 0.850 & 0.850 & +0.000\\
93 & 455 & 0.929 & 0.946 & +0.018 & 10278 & 67 & 0.917 & 0.917 & +0.000 & 12690 & 51 & 1.000 & 1.000 & +0.000\\
100 & 1,753 & 0.943 & 0.965 & +0.022 & 10280 & 7,401 & 0.990 & 0.994 & +0.003 & 12694 & 608 & 1.000 & 1.000 & +0.000\\
103 & 1,094 & 0.982 & 1.000 & +0.018 & 10378 & 982 & 0.976 & 0.979 & +0.003 & 12703 & 315 & 1.000 & 1.000 & +0.000\\
104 & 476 & 0.989 & 0.989 & +0.000 & 10434 & 809 & 0.964 & 0.943 & -0.021 & 12725 & 907 & 1.000 & 0.998 & -0.002\\
105 & 1,365 & 1.000 & 1.000 & +0.000 & 10475 & 258 & 0.847 & 0.903 & +0.056 & 12883 & 149 & 1.000 & 1.000 & +0.000\\
106 & 1,204 & 1.000 & 1.000 & +0.000 & 10495 & 968 & 0.990 & 1.000 & +0.010 & 12896 & 995 & 0.990 & 1.000 & +0.010\\
107 & 19,682 & 0.993 & 0.991 & -0.002 & 10498 & 1,950 & 0.961 & 0.968 & +0.007 & 12909 & 308 & 0.839 & 0.871 & +0.032\\
108 & 1,479 & 1.000 & 1.000 & +0.000 & 10532 & 460 & 1.000 & 1.000 & +0.000 & 12910 & 75 & 0.667 & 0.667 & +0.000\\
114 & 11,621 & 0.973 & 0.979 & +0.006 & 10544 & 690 & 0.940 & 0.956 & +0.016 & 12952 & 5,690 & 0.968 & 0.946 & -0.023\\
116 & 657 & 1.000 & 1.000 & +0.000 & 10548 & 1,123 & 0.812 & 0.821 & +0.009 & 12967 & 175 & 1.000 & 1.000 & +0.000\\
117 & 145 & 1.000 & 1.000 & +0.000 & 10627 & 6,378 & 0.941 & 0.969 & +0.028 & 12968 & 5,754 & 0.953 & 0.952 & -0.002\\
118 & 1,239 & 0.857 & 0.910 & +0.053 & 10635 & 536 & 0.936 & 0.957 & +0.021 & 13000 & 153 & 0.949 & 0.949 & +0.000\\
121 & 2,624 & 0.968 & 0.979 & +0.011 & 10702 & 268 & 0.853 & 0.920 & +0.067 & 13001 & 387 & 0.970 & 0.974 & +0.004\\
124 & 2,019 & 0.966 & 0.974 & +0.008 & 10839 & 749 & 1.000 & 1.000 & +0.000 & 13004 & 328 & 0.730 & 0.972 & +0.242\\
125 & 1,457 & 0.322 & 0.662 & +0.340 & 10911 & 304 & 0.952 & 1.000 & +0.048 & 19904 & 1,449 & 0.958 & 0.959 & +0.001\\
127 & 865 & 0.988 & 0.988 & +0.000 & 10918 & 162 & 0.867 & 0.933 & +0.067 & 19905 & 3,709 & 0.964 & 0.972 & +0.009\\
129 & 8,150 & 0.950 & 0.973 & +0.023 & 10927 & 601 & 0.907 & 0.911 & +0.004 & 20073 & 36 & 0.750 & 0.750 & +0.000\\
130 & 22,153 & 0.989 & 0.969 & -0.019 & 10938 & 303 & 0.867 & 0.967 & +0.100 & 20130 & 86 & 1.000 & 1.000 & +0.000\\
133 & 509 & 0.933 & 0.940 & +0.006 & 10980 & 389 & 0.944 & 1.000 & +0.056 & 20131 & 829 & 0.867 & 0.855 & -0.012\\
134 & 820 & 0.971 & 0.917 & -0.054 & 10982 & 1,027 & 0.971 & 0.961 & -0.010 & 20151 & 584 & 1.000 & 1.000 & +0.000\\
136 & 5,208 & 0.982 & 0.984 & +0.002 & 11016 & 285 & 1.000 & 1.000 & +0.000 & 20154 & 493 & 1.000 & 1.000 & +0.000\\
137 & 9,548 & 0.964 & 0.971 & +0.007 & 11024 & 528 & 0.925 & 0.962 & +0.038 & 20174 & 3,227 & 0.978 & 0.984 & +0.006\\
138 & 3,774 & 0.954 & 0.973 & +0.019 & 11042 & 322 & 1.000 & 1.000 & +0.000 & 30043 & 839 & 0.989 & 0.996 & +0.007\\
142 & 217 & 0.955 & 1.000 & +0.045 & 11060 & 1,023 & 0.891 & 0.891 & +0.000 & 100010 & 1,351 & 0.910 & 0.932 & +0.022\\
144 & 427 & 0.964 & 0.976 & +0.012 & 11063 & 420 & 0.940 & 0.957 & +0.017 & 100138 & 262 & 1.000 & 1.000 & +0.000\\
155 & 1,367 & 0.896 & 0.907 & +0.011 & 11084 & 706 & 0.914 & 0.943 & +0.029 & 100166 & 731 & 0.968 & 0.970 & +0.002\\
165 & 2,225 & 0.994 & 0.995 & +0.002 & 11085 & 856 & 0.956 & 0.962 & +0.006 & 100436 & 207 & 0.847 & 0.882 & +0.035\\
168 & 964 & 1.000 & 1.000 & +0.000 & 11092 & 310 & 1.000 & 1.000 & +0.000 & 100468 & 1,491 & 0.379 & 0.362 & -0.017\\
175 & 861 & 1.000 & 1.000 & +0.000 & 11110 & 301 & 0.967 & 1.000 & +0.033 & 100855 & 407 & 1.000 & 1.000 & +0.000\\
176 & 12,589 & 0.809 & 0.823 & +0.014 & 11140 & 940 & 1.000 & 1.000 & +0.000 & 100990 & 54 & 1.000 & 1.000 & +0.000\\
184 & 175 & 1.000 & 1.000 & +0.000 & 11149 & 232 & 0.952 & 0.984 & +0.032 & 101365 & 66 & 0.688 & 0.688 & +0.000\\
185 & 3,009 & 0.985 & 0.997 & +0.012 & 11154 & 209 & 0.975 & 0.975 & +0.000 & 102439 & 448 & 0.978 & 0.978 & +0.000\\
193 & 698 & 0.797 & 0.862 & +0.065 & 11156 & 1,271 & 0.836 & 0.891 & +0.055 & 102728 & 164 & 1.000 & 1.000 & +0.000\\
194 & 7,329 & 0.976 & 0.984 & +0.008 & 11261 & 191 & 1.000 & 1.000 & +0.000 & 103154 & 1,978 & 0.716 & 0.721 & +0.005\\
216 & 498 & 1.000 & 1.000 & +0.000 & 11272 & 1,581 & 0.983 & 0.987 & +0.004 & 103454 & 994 & 1.000 & 1.000 & +0.000\\
218 & 463 & 1.000 & 0.978 & -0.022 & 11280 & 683 & 0.996 & 1.000 & +0.004 & 103767 & 193 & 0.952 & 0.940 & -0.012\\
219 & 1,072 & 0.911 & 0.933 & +0.022 & 11287 & 980 & 0.985 & 0.998 & +0.013 & 104320 & 5,198 & 0.883 & 0.955 & +0.072\\
222 & 98 & 0.900 & 0.900 & +0.000 & 11290 & 2,054 & 0.994 & 0.996 & +0.002 & 104499 & 119 & 0.955 & 1.000 & +0.045\\
226 & 445 & 0.982 & 1.000 & +0.018 & 11291 & 230 & 1.000 & 1.000 & +0.000 & 104573 & 729 & 0.882 & 0.924 & +0.042\\
227 & 868 & 0.983 & 0.988 & +0.006 & 11336 & 709 & 0.940 & 0.955 & +0.015 & 105440 & 59 & 0.833 & 0.833 & +0.000\\
235 & 312 & 0.958 & 1.000 & +0.042 & 11362 & 4,631 & 0.937 & 0.860 & -0.077 & 105674 & 182 & 0.600 & 0.800 & +0.200\\
236 & 181 & 0.757 & 0.814 & +0.057 & 11400 & 1,239 & 0.996 & 0.996 & +0.000 & 109939 & 644 & 1.000 & 1.000 & +0.000\\
243 & 927 & 0.974 & 0.989 & +0.015 & 11402 & 1,084 & 0.969 & 0.985 & +0.015 &  &  &  &  & \\
246 & 225 & 0.900 & 0.850 & -0.050 & 11440 & 5,811 & 0.910 & 0.928 & +0.017 &  &  &  &  & \\
\end{longtable}

\normalsize
\end{landscape}

\end{document}
